\section{Теоретическая часть}
\subsection{Формулы для вычисления величин}
Пусть \( \vec{X} = (X_1, \ldots, X_n) \) --- случайная выборка из генеральной
совокупности \( X \), а \( \vec{x} = (x_1, \ldots, x_n) \) --- реализация этой выборки.
Расположим элементы выборки \( \vec{x} \) в порядке неубывания, \( i \)-тый элемент 
полученной последовательности обозначим \( x_{(i)} \), \( i = \overline{1, n} \): 
\( x_{(1)} \leq \dots \leq x_{(n)} \).

Тогда максимальное значение выборки вычисляется по формуле~\ref{eq:s-max},
\begin{equation} \label{eq:s-max}
    M_{max} = x_{(n)} = \max\{x_1, \ldots, x_n\},
\end{equation}

минимальное значение выборки вычисляется по формуле~\ref{eq:s-min},
\begin{equation} \label{eq:s-min}
    M_{min} = x_{(1)} = \min\{x_1, \ldots, x_n\},
\end{equation}

размах вычисляется по формуле~\ref{eq:s-ampl},
\begin{equation} \label{eq:s-ampl}
    R = M_{max} - M_{min} = x_{(n)} - x_{(1)},
\end{equation}

выборочное среднее вычисляется по формуле~\ref{eq:mean},
\begin{equation} \label{eq:mean}
    \hat{\mu}(\vec{x}) = \overline{x} = \frac{1}{n}\sum_{i=1}^{n}x_i,
\end{equation}

исправленная выборочная дисперсия вычисляется по формуле~\ref{eq:var},
\begin{equation} \label{eq:var}
    S^2(\vec{x}) = \frac{1}{n-1}\sum_{i=1}^{n}(x_i - \bar{x} )^2.
\end{equation}

\subsection{Определение эмпирической функции плотности и гистограммы}
Пусть
\begin{itemize}
    \item \( \vec{X} = (X_1, \ldots, X_n) \) --- случайная выборка из распределения случайной
        величины \( X \), а \( \vec{x}_n = (x_1, \ldots, x_n) \) --- реализация этой выборки;
    \item отрезок \( [x_{(1)}, \enspace x_{(n)} ] \) разбит на \( m \in \mathbb{N} \) непересекающихся полуинтервалов
        \( J_1, \ldots, J_m \)  длины \( \Delta = \frac{(x_{(n)} - x_{(1)})}{m} \), где
        \begin{equation*}
            J_i = [x_{(1)} + \Delta \cdot (i - 1); \enspace x_{(1)} + \Delta \cdot i ), i = \overline{1, (m - 1)},
        \end{equation*} 
        \begin{equation*}
            J_m = [x_{(n)} - \Delta; \enspace x_{(n)}];
        \end{equation*}
    \item \( n_i \) --- число наблюдений, попавших в полуинтервал \( J_i \), \( i = \overline{1, m} \);
\end{itemize}
тогда эмпирической функцией плотности называется функция
\begin{equation*}
    f_n(x) =
     \begin{cases}
        \frac{n_i}{\Delta \cdot n}, x_i \in J_i, \\
        0, \text{иначе}. 
    \end{cases}
\end{equation*}

График эмпирической функции плотности называетя гистограммой.

\subsection{Определение эмпирической функции распределения}
Пусть
\begin{itemize}
    \item \( \vec{X} = (X_1, \ldots, X_n) \) --- случайная выборка из распределения случайной
        величины \( X \), а \( \vec{x}_n = (x_1, \ldots, x_n) \) --- реализация этой выборки;
    \item \( h(x, \vec{x}) \) --- количество элементов выборки \( \vec{x} \), 
        которые имеют значение, меньшее \( x \),
\end{itemize}
тогда функция \( F_n(x) = \frac{h(x, \vec{x})}{n} \) называется эмпирической функцией
распределения, отвечающей выборке \( \vec{x} \). 